\documentclass[11pt]{article}
\usepackage[T1]{fontenc}
\usepackage[utf8]{inputenc}
\usepackage{lmodern,microtype}
\usepackage[margin=1in]{geometry}
\usepackage{amsmath,amssymb,amsthm,mathtools}
\usepackage{array,booktabs,longtable,enumitem}
\usepackage{xcolor}
\usepackage{url}
\usepackage{hyperref}
\usepackage{fancyhdr}
\definecolor{inkblue}{RGB}{27,57,88}
\hypersetup{colorlinks=true,linkcolor=inkblue,citecolor=inkblue,urlcolor=inkblue,
 pdftitle={From finite presentations to explicit witnesses in iPRA},
 pdfauthor={Research note prepared for Vladimir Reshetnikov}}
\newcommand{\ipra}{\ensuremath{\mathrm{iPRA}}}
\newcommand{\iPRA}{\ipra}
\newcommand{\N}{\mathbb N}
\newcommand{\Z}{\mathbb Z}
\newcommand{\Q}{\mathbb Q}
\newcommand{\PR}{\mathrm{PR}}
\newcommand{\CertCheck}{\operatorname{Check}}
\newcommand{\eval}{\operatorname{eval}}
\newcommand{\iter}{\operatorname{it}}
\newcommand{\len}{\operatorname{len}}
\newcommand{\code}[1]{\texttt{#1}}
\newcolumntype{P}[1]{>{\raggedright\arraybackslash}p{#1}}
\newtheorem{theorem}{Theorem}[section]
\newtheorem{lemma}[theorem]{Lemma}
\newtheorem{proposition}[theorem]{Proposition}
\newtheorem{corollary}[theorem]{Corollary}
\theoremstyle{definition}
\newtheorem{definition}[theorem]{Definition}
\newtheorem{example}[theorem]{Example}
\theoremstyle{remark}
\newtheorem{remark}[theorem]{Remark}
\setlist{itemsep=3pt,topsep=5pt}
\setlength{\emergencystretch}{2em}
\pagestyle{fancy}
\fancyhf{}
\fancyhead[L]{\small Finite presentations and witnesses in \ipra}
\fancyhead[R]{\small Research note}
\fancyfoot[C]{\thepage}
\setlength{\headheight}{14pt}
\title{\textbf{From finite presentations to explicit witnesses in \ipra}\\[6pt]
\large Proof-preserving simulation, finite observable modules,\\
and complete rational-point enumeration}
\author{Research note prepared for Vladimir Reshetnikov}
\date{15 November 2026}
\begin{document}
\maketitle
\begin{abstract}
We develop a reusable route from finitely presented mathematical data
to explicitly witnessed theorems in intuitionistic primitive recursive
arithmetic. The first component is a syntactic simulation theorem:
an actual arithmetic derivation of a local simulation law produces
an actual derivation for every finite iteration. Its equality case
transports iterative constructions through quotients, evaluations,+and truncations. Both transformations are implemented and checked in Lean.

The second component extracts more information from an existing
finite-algebra interface in ProveIt. A rational unit-and-multiplication
certificate supplies explicit coordinate denominators and absolute
bounds, and therefore a complete finite enumeration of rational zeros.
The proof uses primitive integer eigenvectors and a column-sum bound;
it needs neither a basis nor a determinant. We give a certificate
constructor from separated annihilating polynomials, extensions to
open and projective charts, and a complementary complete enumeration
procedure for plane sections with a nonzero polynomial norm.

The distinction between a supplied global cover and covers depending
on a queried finite tuple is made explicit. The results strengthen
several finite algebraic interfaces in the repository, but do not
establish general Siegel, Faltings, Mordell--Weil, or the unresolved
analytic estimates. The archive contains full human proofs, exact
reference programs, compiled Lean cores, complete definitions for the
remaining Lean targets, and a separate route to object-language
arithmetic derivations.
\end{abstract}
\tableofcontents
\newpage
\section{Scope, source audit, and the contribution}
\label{sec:scope}

The purpose of this note is to make mathematical work in \ipra\ more
composable while actually improving one of the repository's interfaces.
The supplied repository is substantial: the inspected snapshot contains
sixty theorem packets, including nine earlier frameworks. Their index
explicitly distinguishes the checked ConciseIPRA development from packets
filed as delivered without a repository-wide verification pass
\cite{RepoIndex,PacketIndex}. Repeating their finite-observation vocabulary
would not resolve an additional mathematical obligation.

We therefore pin the audit to commit
\begin{center}\small
\path{b0f6f535ffc16324683fef2c263260a88ae03755}
\end{center}
and concentrate on two missing connections. The first connects a local
proof about a change of representation to a derivation about an entire
finite computation. The second connects an algebraic certificate of
finite dimensionality to positive arithmetic witnesses: coordinate bounds
and complete rational lists.

\subsection{Results, in their exact logical forms}

The main logical result has a deliberately syntactic premise. Let
$f,g,h$ be fixed unary PR function symbols. From a supplied derivation
in \ipra\ of
\[
                      h(f(x))=g(h(x))
\]
we construct a derivation of
\[
       \forall n\;h(\iter_f(n,x))=\iter_g(n,h(x)).
\]
The proof uses the displayed equality itself as a quantifier-free
induction invariant. The companion theorem transports an arbitrary
quantifier-free binary relation preserved by a pair of steps.
These are checked Lean theorems with \code{Prf PRAAx} conclusions;
they are not a rule promoting a theorem of the host logic to an
arithmetic axiom. A concrete repeated-squaring corollary connects this
infrastructure to polynomial evaluation, modular computations and
truncated series.

The main algebraic result starts with polynomial identities, all checked
by finite coefficient arithmetic. Suppose $F_1,\ldots,F_s$ are rational
polynomials in $r$ variables and we are supplied a finite family $b_i$,
a unit combination, and multiplication matrices $M_k$. At a common
zero $p$, the evaluation row $v=(b_i(p))$ is nonzero and satisfies
\[
                           vM_k=p_kv.
\]
Choose a positive integer $\Delta_k$ clearing the entries of $M_k$,
and put
\[
 A_k=\Delta_kM_k,\qquad
 R_k=\max_j\sum_i |(A_k)_{ij}|.
\]
We prove
\[
             p_k\in\{z/\Delta_k:-R_k\le z\le R_k\}.
\]
Consequently exhaustive exact filtering of at most
$\prod_k(2R_k+1)$ tuples returns \emph{all} rational common zeros.
No height bound is an extra input to this theorem: the accepted matrices
compute it. An abstract assertion that there are finitely many rational
zeros is not an adequate input.

The norm result addresses a second existing interface. For a monic plane
equation $F(X,Y)$ and a supplied section $G(X,Y)$, compute the norm of
$G$ in the free $\Q[X]$-module $\Q[X,Y]/(F)$. If that polynomial
norm is nonzero, bounded rational-root search gives a complete list of
the common rational zeros. Thus a known section can be located, not
merely counted. A global supplied finite covering list of sections can
then be enumerated by finite union. A tuple-dependent covering assertion
does not provide such a list.

\subsection{Relationship with the existing frameworks}

\begin{longtable}{P{.23\linewidth}P{.34\linewidth}P{.33\linewidth}}
\toprule
Inspected interface & Existing content & Additional result here\\
\midrule\endhead
Finite contracts and observations & Booleanized bounded invariants,
precision demands and guarded constructions. & A compiled transformation
of local object derivations into arbitrary finite simulations.\\[4pt]
Coherent methods & Finite uniqueness, exact restriction maps, filtered
lifting and object-level contract composition. & A general iteration
transport theorem and its relational version.\\[4pt]
Finite algebra certificates & A size-$D$ spanning certificate bounds
the length of a list of distinct rational zeros by $D$. & The same
action matrices give denominator divisibility, coordinate bounds and
complete rational enumeration.\\[4pt]
Finite norm certificates & Polynomial norms control the number of
points in a supplied algebraic section. & A nonzero norm gives a
complete section list by two finite rational-root searches.\\[4pt]
Finite covers in the Siegel route & Tuple-local proper-subspace covers
give a cardinality implication. & A precise stronger input, an actual
global list or a coefficient bound for the sections, gives enumeration.\\
\bottomrule
\end{longtable}

The direct finite-algebra predecessor is
\path{foundations/ipra_embedding_framework/paper/05_diophantine.tex},
especially its subsection on finite algebra certificates
\cite{Embedding}. The most recent norm and cover interfaces occur in
\path{ipra_finite_demand} and \path{ipra_coherent_methods_2026-11-15}
\cite{Demand,Coherent}. The earlier
\path{ipra_finite_observation/tex/sections.tex} explicitly explains
why cardinality and finite-tuple certificates alone yield no global
height cutoff \cite{Observation}. We use those distinctions and
strengthen the data-rich algebraic cases; we do not contradict the
general obstruction.

Multiplication matrices and effective solving of zero-dimensional
polynomial systems are established methods in computer algebra.
Rouillier's rational univariate representation is an important example
\cite{Rouillier}. This paper does not claim that these broad ideas are
new. The contribution relative to the inspected repository is the
small certificate-to-grid theorem, its weak-arithmetic proof,
the explicit constructors and scope distinctions, and the checked
object-language transport component.

\subsection{Three levels of verification}

We distinguish the following assertions throughout.
\begin{enumerate}
\item A \emph{human \ipra\ argument} describes fixed PR algorithms and
proves them using finite identities and explicit quantifier-free
induction invariants. The new algebraic and norm results have such
arguments here.
\item A \emph{host Lean theorem} proves an assertion about natural
numbers, integers, lists or rational arithmetic in Lean. The compiled
integer-eigenvalue lemmas and finite-filter reduction have this status.
\item An \emph{object Lean theorem} has a conclusion in the inductively
defined relation \code{Prf PRAAx}. The compiled simulation, fusion and
repeated-squaring transport theorems have this stronger, syntactic status.
\end{enumerate}
The full rational enumeration theorem has not yet been reified as a
\code{Prf PRAAx} derivation. Its Lean file defines the finite data and
programs and states the remaining obligations as named propositions,
without axioms or placeholder proofs. Section~\ref{sec:formalization}
gives the dependency order and the precise final arithmetic formula.

\subsection{What this does not resolve}

The reported difficulty in general Mordell--Weil is producing and
certifying the covering representatives, not running a descent after
they have been supplied \cite{MordellWeil}. The general Siegel and
Faltings packets do not supply the global effective data required for
their original requests \cite{Siegel,Faltings}. The analytic
prime-distribution packet leaves substantial uniform estimates open
\cite{PrimeDistribution}. The present framework supplies neither those
estimates nor a hidden choice procedure. It does supply new, fully
proved algebraic subroutines and makes a completed formal arithmetic
interface closer for them.

The audit is targeted. We inspected the repository index, all nine
framework descriptions, their relevant sections, the directly affected
arithmetic-geometry reports, and the actual ConciseIPRA definitions.
We did not independently verify every theorem in all sixty packets.

\section{A framework based on finite presentations and proof transport}
\label{sec:framework}

\subsection{The arithmetic base and the meaning of an embedding}

Our base theory is first-order intuitionistic logic with equality,
one function symbol for each primitive recursive definition, their
defining equations, successor injectivity and nonzeroness, and induction
for quantifier-free formulas, with parameters. This is the presentation
implemented by the repository's \code{PR}, \code{PTm}, \code{PRAAx}
and \code{Prf} types \cite{PRSource,DeductionSource}. We do not add
induction for an arbitrary universally quantified assertion, excluded
middle for arbitrary propositions, choice, or an axiom of convergence.

The term ``embedding'' here means a specified translation of the finite
content of a mathematical statement into arithmetic. It does not mean
that every theorem of a stronger ambient theory becomes a theorem of
\ipra. The translation records the representation, the observations
being requested, a concrete algorithm, and a derivation of its stated
finite contract. Distinct translations can preserve different information.
A bound on the number of points, a list of all points, a rational
approximation and a coherent family of approximations are distinct outputs.

Use any explicitly fixed PR pairing and list coding. For definiteness
one may pair $a,b$ by $(a+b)(a+b+1)/2+b$, and encode a nonempty list
by one plus the pair of its head and tail. Projections are bounded
searches; a nonzero tail code decreases. All list folds therefore have
PR definitions with fuel at most the input code. More efficient list
implementations are extensionally interchangeable after their simulation
law has been proved. Encode integers by a sign bit and magnitude, and
rationals by a signed numerator and a positive denominator, normalized
by gcd; zero has denominator one. A malformed code has a designated
default and a separate validity flag. A theorem never relies on the
default satisfying the specification of a valid mathematical object.

Rational equality, order and arithmetic reduce to signed integer
arithmetic and positive-denominator cross multiplication. The Euclidean
algorithm can be run for $a+b+1$ steps on nonnegative inputs $a,b$,
freezing the state when the second component becomes zero. Each active
step strictly decreases a nonnegative component. Its invariant tracks
the common divisors and, for extended gcd, two linear-combination
identities. Thus division, gcd, lcm and rational normalization do not
require unbounded minimization. Sparse multivariate polynomials are
finite lists of coefficient/exponent-vector pairs; evaluation is a fixed
fold using finite powers. These concrete data operations will also be
defined in the Lean target file.

\subsection{Bounded formulas and allowed induction}

For a fixed PR Boolean predicate $P(a,i)$, put
\[
 A(a,0)=1,\qquad A(a,t+1)=A(a,t)\mathbin{\wedge}P(a,t).
\]
Then a bounded assertion $\forall i<t\,P(a,i)=1$ is represented by
the atomic equality $A(a,t)=1$. A corresponding disjunction fold
represents bounded existence. The familiar unfolding equivalences
are proved with the queried index free. For the converse under a fixed
bound $B$, use the guarded quantifier-free invariant
$t\le B\to A(a,t)=1$, with the bounded-universal premise kept
outside the induction formula. For nested bounded
quantifiers the construction is repeated, from the inside outward.
This is a schema for each fixed formula and its PR predicates.

Finite-list correctness assertions, such as ``each processed coefficient
agrees'' or ``all entries satisfy this divisibility equation'', can
therefore be used as quantifier-free invariants. When a proof concerns
an arbitrary solution $p$, keep $p$ as a free parameter during induction.
Do not put ``for every rational point'' inside the induction invariant.
This simple distinction matters in the enumeration proof: the finite
matrix calculation is performed with one arbitrary point fixed, and
universal generalization occurs afterwards.

There is no total PR interpreter for arbitrary PR program codes. If
one had $E(e,x)$ and every unary PR function were represented, the
PR function $d(x)=E(x,x)+1$ would have a code $e$ and give
$d(e)=d(e)+1$. Our constructions avoid that contradiction. The simulation
theorems are schemas for specified PR syntax. The polynomial checker
is uniform over arbitrary polynomial \emph{data}, because it interprets
one fixed polynomial grammar with one fixed PR evaluation algorithm.
The two forms of uniformity are compatible and must be stated separately.

\subsection{Finite presentations and observations}

A finite presentation has a validity predicate $V(a)$ and, when needed,
a decidable equivalence relation $a\sim b$ on valid codes. A restricted
operation has a total PR implementation and a separate guard under which
its mathematical law holds. Examples include division with a nonzero
denominator, a matrix inverse with a verified inverse table, and a
localization coordinate with a nonzero polynomial value.

An observation is a fixed PR function or relation on the representation.
It may be an exact polynomial coefficient, a matrix entry, a residue
modulo $q^n$, a rational number with an error bound, or membership in a
finite combinatorial object. A mathematical proof is organized by the
observations it actually uses. Infinite terminology is explanatory until
the corresponding arithmetic observations and their precision demands
have been specified.

For example, a metric name $a(k)\in\Q$ can be required to satisfy
$|a(k)-a(\ell)|\le2^{-k}+2^{-\ell}$ at the finitely many indices
needed in a calculation. A formal series is observed by exact coefficient
truncation. A residue tower is observed modulo a specified power. These
are different contracts. There is no rule converting a formal series
into an analytic function without a tail estimate, or a local residue
tower into a rational point without a global argument.

\subsection{Syntactic iteration fusion}

For a specified unary PR operation $f$, define
\[
 \iter_f(0,x)=x,\qquad
 \iter_f(n+1,x)=f(\iter_f(n,x)).
\]
This definition is itself primitive recursion. In the repository syntax,
the recursion counter comes first and the previous result is the second
input to the step. The Lean construction chooses exactly those projections.

\begin{theorem}[Proof-preserving fusion]\label{thm:fusion}
Let $f,g,h$ be specified unary PR function symbols. If an actual
\ipra\ derivation of $h(f(x))=g(h(x))$ is supplied, a finite
syntactic construction produces an actual \ipra\ derivation of
\[
 \forall n\;h(\iter_f(n,x))=\iter_g(n,h(x)).
 \tag{F}\label{eq:fusion}
\]
\end{theorem}

\begin{proof}
Use the equality in \eqref{eq:fusion} as the quantifier-free invariant
$I(n,x)$. At zero, the two recursion equations give $h(x)=h(x)$.
For the successor step, start with the induction hypothesis and compute
inside the object equality calculus:
\begin{align*}
 h(\iter_f(n+1,x))
 &=h(f(\iter_f(n,x)))\\
 &=g(h(\iter_f(n,x)))\\
 &=g(\iter_g(n,h(x)))\\
 &=\iter_g(n+1,h(x)).
\end{align*}
The second equality is the supplied local derivation after substituting
$\iter_f(n,x)$ for its variable; the third is equality congruence
applied to the induction hypothesis. The first and last use defining
equations. Apply the object rule for quantifier-free induction, keeping
$x$ free. The operations used on derivations are substitution,
weakening, congruence, transitivity and the induction constructor.
Every one is an existing operation on \code{Prf}; no semantic reflection
step occurs. \code{ObjectFusion.fusion} implements this proof.
\end{proof}

The theorem also covers algorithms with numerical parameters. Put the
parameter in the state and use the unary step $(a,x)\mapsto(a,F(a,x))$.
A representation map can preserve that parameter. This is a fixed
PR pair operation, not a family of programs selected by a universal
interpreter. A step that depends on its iteration index is handled by
also storing and incrementing the counter in the state.

\subsection{Relations, guards and simulation}

Equality is often too restrictive. Quotient representatives may differ,
and rational approximations carry inequalities instead of exact equalities.
The relational theorem gives the appropriate common statement.

\begin{theorem}[Quantifier-free simulation]\label{thm:simulation}
Let $R(x,y)$ be a specified quantifier-free arithmetic formula and let
$f,g$ be specified unary PR steps. From a derivation
\[
                      R(x,y)\ \vdash\ R(f(x),g(y))
\]
one can construct a derivation
\[
       R(x,y)\ \vdash\ \forall n\,
                      R(\iter_f(n,x),\iter_g(n,y)).
 \tag{S}\label{eq:simulation}
\]
\end{theorem}

\begin{proof}
Substitution into a quantifier-free formula preserves its quantifier-free
status, by structural inspection of the formula. Hence
$R(\iter_f(n,x),\iter_g(n,y))$ is an allowed induction invariant.
At zero, the hypothesis $R(x,y)$ and the two defining equations give
the base case. Substitute the two current iterates for $x,y$ in the
supplied local derivation. Its premise is the induction hypothesis;
its conclusion becomes the next invariant after applying the successor
equations. Quantifier-free induction proves \eqref{eq:simulation}.
The formal proof keeps the shifted initial premise in the context under
the counter binder. \code{ObjectSimulation.simulation} checks these
substitution and variable-scope details explicitly.
\end{proof}

Guards should be part of $R$. For example,
$R(x,y)$ may assert that both codes are valid and that their evaluations
agree, or that both denominators are positive and an integer-cleared
error inequality holds. A local lemma valid only on well-formed inputs
then preserves the entire invariant. This prevents a later iteration
from invoking a law outside its domain.

\begin{corollary}[Transport through repeated squaring]
\label{cor:squaring}
Suppose a specified map $h$ between two coded multiplication systems
satisfies the object derivation
$h(a\cdot b)=h(a)\star h(b)$.
Define $s_0=x$, $s_{n+1}=s_n\cdot s_n$, and
$t_0=h(x)$, $t_{n+1}=t_n\star t_n$. Then \ipra\ proves
$h(s_n)=t_n$ for every $n$.
\end{corollary}
\begin{proof}
Substitute $a=b=x$ into the local multiplication derivation to obtain
the premise of Theorem~\ref{thm:fusion} for the two squaring steps.
Apply that theorem. No associativity is needed for this particular
statement; an identification with $x^{2^n}$ may be proved separately
in an associative system. The supplied Lean corollary performs the
substitution and invokes the checked fusion builder.
\end{proof}

The same idea applies to a polynomial expression evaluated by an
instruction list: store its program counter, stack and fixed parameters
in the state. A local stack-operation simulation yields correctness of
the whole finite evaluator. This is an interpretation of a fixed finite
expression grammar, whose operations are specified; it is not an
interpreter for all PR functions.

\subsection{Exact witnesses from a finite observable image}

The algebraic theorem in the next section supplies a finite grid
containing all possible rational values of selected observables. Here
is the general step that turns such a grid into object witnesses.

\begin{proposition}[Finite reconstruction]\label{prop:reconstruction}
Let $V(p)$ be a decidable predicate on finite codes, let
$o(p)$ be a fixed PR observation, and let $G$ be an explicitly computed
finite list. Suppose the following statements have \ipra\ proofs:
\begin{enumerate}
\item $V(p)$ implies $o(p)\in G$;
\item for each $q\in G$, a fixed PR list producer $B(q)$ contains
every valid $p$ with $o(p)=q$;
\item membership in $B(q)$ and all validity and observation tests are
finite decidable tests.
\end{enumerate}
Then filter $\bigcup_{q\in G}B(q)$ by $V$ and remove duplicates.
The resulting PR list contains exactly the valid objects.
\end{proposition}
\begin{proof}
A retained object satisfies $V$ by its filter test. Conversely, fix an
arbitrary valid $p$. The first premise supplies its observation among
the finite $G$; the second supplies $p$ in that observation's list.
Therefore $p$ is visited and retained. A finite deduplication pass
preserves membership. Its invariant equates membership in the output
prefix with membership in the processed input prefix; bounded folds
make that invariant quantifier-free. This proves both directions.
\end{proof}

When the observation is the full rational coordinate tuple, $B(q)$
is the singleton $[q]$. This is the important algebraic specialization.
When an observation is constant on an infinite family, there is no
finite reconstruction list in general. The framework exposes the
missing premise instead of concealing it in the word ``representation''.

\section{Finite observable modules produce complete rational point lists}
\label{sec:finite-observable-modules}

The earlier finite-algebra interface proves that a multiplication-closed
spanning certificate of size $D$ bounds the number of rational zeros by
$D$.  There is more information in the same certificate.  Its matrices
also bound every coordinate and every reduced denominator.  The resulting
finite search returns the complete rational zero set, with no externally
supplied height bound.  This is a strengthening of that interface, not a
claim that an arbitrary cardinality bound determines a search cutoff.
The additional information is the supplied multiplication action.

The central observation needs neither a determinant nor a linearly
independent basis.  A rational eigenvalue of an integer matrix is an
integer: after scaling an eigenvector to a primitive integer vector,
the denominator would otherwise divide all its coordinates.  We give
the finite proof because it both simplifies the argument and identifies
exactly the arithmetic needed in iPRA.

\subsection{Primitive integer eigenvectors}

All finite vectors and matrices in this section have at least one
coordinate.  A vector $w\in\mathbb Z^D$ is \emph{primitive} when the
greatest common divisor of $|w_1|,\ldots,|w_D|$ is one.  Rational
numbers are represented by reduced numerator--positive-denominator
pairs.  Matrix multiplication uses row vectors: $(wA)_j=
\sum_{i=1}^D w_i A_{ij}$.

\begin{lemma}[Primitive scaling]
\label{lem:primitive-scaling}
From a supplied nonzero vector $v\in\mathbb Q^D$ one can primitively
recursively construct a nonzero rational $c$ and a primitive integer
vector $w$ such that $w=cv$.  If $vM=\lambda v$, then
$wM=\lambda w$.
\end{lemma}

\begin{proof}
Write the entries of $v$ with positive denominators, and let $L$ be
their product.  Then $u=Lv$ is an integer vector and is nonzero.
Compute $g=\gcd(|u_1|,\ldots,|u_D|)$ by the Euclidean algorithm,
folded along the finite list.  Since $u$ is nonzero, $g>0$.
Every $u_i$ is divisible by $g$, so $w_i=u_i/g$ is an integer and
the gcd of its absolute coordinates is one.  Take $c=L/g$.
Multiplying the displayed eigenvector equation by $c$ gives the final
assertion.  Products, gcd, exact integer division with its failure
branch, and finite folds are primitive recursive.  Extended gcd
additionally produces integers $e_i$ satisfying
$\sum_i e_iw_i=1$, if that form of primitivity is preferred.
\end{proof}

\begin{lemma}[Integer eigenvalue and an explicit bound]
\label{lem:integer-eigenvalue}
Let $A\in\mathbb Z^{D\times D}$, let $w\in\mathbb Z^D$ be
primitive, and let $\lambda\in\mathbb Q$.  If
$wA=\lambda w$, then $\lambda\in\mathbb Z$.  Moreover, with
\[
 R(A)=\max_{1\le j\le D}\sum_{i=1}^D|A_{ij}|,
\]
one has $|\lambda|\le R(A)$.
\end{lemma}

\begin{proof}
Write $\lambda=a/b$ in reduced form, where $b>0$ and
$\gcd(|a|,b)=1$.  For each coordinate $j$, multiplication by $b$
in the eigenvector equation gives
\[
 a w_j=b\sum_iw_iA_{ij}.
\]
Thus $b$ divides $a w_j$, and coprimality gives $b\mid w_j$.
One elementary justification uses a Bezout identity $ra+sb=1$:
multiply it by $w_j$ and substitute the preceding divisibility.
Since $w$ is primitive, there are integers $e_j$ with
$\sum_j e_jw_j=1$.  Therefore $b\mid1$.  Positivity of $b$
forces $b=1$.

For the bound, choose the first index $j$ attaining
$W=\max_i|w_i|$.  The finite maximum and its first maximizing
index are obtained by a bounded scan.  Primitivity implies $W>0$.
The $j$-th coordinate gives
\[
 |\lambda|W
   =\left|\sum_iw_iA_{ij}\right|
   \le\sum_i|w_i||A_{ij}|
   \le W\sum_i|A_{ij}|
   \le WR(A).
\]
Cancel the positive integer $W$.  No matrix inverse, determinant,
characteristic polynomial, or spectral theorem was used.
\end{proof}

\begin{corollary}[A rational eigenvalue lies on a finite grid]
\label{cor:eigenvalue-grid}
Let $M\in\mathbb Q^{D\times D}$ and let $\Delta>0$ be an
integer such that $A=\Delta M$ has integer entries.  If a nonzero
$v\in\mathbb Q^D$ satisfies $vM=xv$, then
\[
 x=\frac z\Delta
 \quad\text{for some }z\in\mathbb Z\text{ with }|z|\le R(A).
\]
In particular, if $x=a/q$ is reduced with $q>0$, then
$q\mid\Delta$, $q\le\Delta$, and $|a|\le R(A)$.
\end{corollary}

\begin{proof}
Primitive scaling gives $wM=xw$ with $w$ primitive, and hence
$wA=(\Delta x)w$.  Lemma~\ref{lem:integer-eigenvalue} says that
$z=\Delta x$ is an integer with the asserted bound.  To reduce
$z/\Delta$, divide numerator and denominator by
$\gcd(|z|,\Delta)$.  This proves the denominator divisibility
and numerator bounds, including $z=0$.
\end{proof}

\subsection{The observable-module principle}

This principle applies before choosing a particular polynomial geometry.
Let $S$ be a class of finitely coded objects with a decidable validity
predicate.  Let $u_1,\ldots,u_r:S\to\mathbb Q$ be the rational
observables whose possible values are wanted.  Fix rational-valued
functions $b_1,\ldots,b_D$, rational numbers $c_1,\ldots,c_D$,
and rational matrices $M_1,\ldots,M_r$.  Suppose there are proved
local contracts, for every valid $p$, of the form
\begin{align}
  1&=\sum_{i=1}^D c_i b_i(p),
      \label{eq:observable-unit}\\
  u_k(p)b_j(p)&=\sum_{i=1}^D (M_k)_{ij}b_i(p)
      &&(1\le k\le r, 1\le j\le D).
      \label{eq:observable-action}
\end{align}
The names ``observable module'' and ``spanning family'' refer only to
these displayed finite identities.  One need not introduce an abstract
module or assert that the functions $b_i$ are linearly independent.

\begin{theorem}[Finite observable image]
\label{thm:finite-observable-image}
Choose positive common denominators $\Delta_k$ for the matrices,
put $A_k=\Delta_kM_k\in\mathbb Z^{D\times D}$, and put
$R_k=R(A_k)$.  Every valid $p$ satisfies
\[
 (u_1(p),\ldots,u_r(p))\in
 \prod_{k=1}^r\left\{\frac z{\Delta_k}:z\in\mathbb Z,
                                      -R_k\le z\le R_k\right\}.
 \label{eq:observable-grid}
\]
Consequently the image has at most $\prod_k(2R_k+1)$ elements.
If all functions and the validity test have specified primitive
recursive implementations and the local contracts have iPRA proofs,
the displayed implication also has an iPRA proof, uniformly in the
finite matrix and denominator data.
\end{theorem}

\begin{proof}
For a fixed valid $p$, let $v=(b_1(p),\ldots,b_D(p))$.
The unit identity implies that $v\ne0$: if all entries vanished,
their indicated linear combination would be zero, contradicting one.
The action identity says precisely that $vM_k=u_k(p)v$.
Apply Corollary~\ref{cor:eigenvalue-grid} separately to every $k$.
The product of the finite coordinate grids contains every image value.
Each grid has exactly $2R_k+1$ entries because $\Delta_k>0$,
and finite tuple enumeration gives the cardinality bound.

For the iPRA assertion, the arbitrary coded input $p$ is a parameter
of the proof.  All remaining quantifiers in the matrix-vector
identities, gcd folds, Bezout checks, and maximum searches range over
supplied finite lists.  Compile them to Boolean folds.  The Euclidean
and list identities just proved then use quantifier-free induction.
There is no quantification over all functions or all possible programs.
\end{proof}

If observables fail to determine $p$, this theorem supplies a finite
\emph{image}, not a list of all input objects.  An infinite set can
have a constant observable.  Complete enumeration therefore needs
actual rational coordinates, or a separately proved finite reconstruction
procedure from the observable values.  This is a useful distinction:
the theorem can also classify finite invariant values of an infinite
family without pretending to enumerate that family.

\subsection{Polynomial certificates and the enumeration theorem}

Let $F_1,\ldots,F_s\in\mathbb Q[X_1,\ldots,X_r]$ be finitely
encoded polynomials.  A \emph{unit-and-action certificate} comprises
polynomials $b_1,\ldots,b_D$, a rational vector $c$, rational
matrices $M_k$, and coefficient-polynomial witnesses for the identities
\begin{align}
 1-\sum_i c_i b_i&=\sum_{\ell=1}^s h_\ell F_\ell,
                \label{eq:module-unit-certificate}\\
 X_kb_j-\sum_i (M_k)_{ij}b_i
       &=\sum_{\ell=1}^s h_{kj\ell}F_\ell.
                \label{eq:module-action-certificate}
\end{align}
The checker expands each side and compares the finitely many
coefficients.  The earlier spanning-certificate interface is the
special case $b_1=1$, $c_1=1$, and $c_i=0$ for $i>1$.
Our relaxed unit condition is occasionally convenient when a generating
family is supplied without a distinguished constant member.

\begin{theorem}[Complete rational enumeration from a finite certificate]
\label{thm:module-enumeration}
From a valid unit-and-action certificate, a fixed primitive recursive
algorithm produces a duplicate-free list $L$ satisfying
\[
  p\in L\quad\Longleftrightarrow\quad
           F_1(p)=\cdots=F_s(p)=0
       \qquad(p\in\mathbb Q^r).
\]
It examines at most $\prod_{k=1}^r(2R_k+1)$ candidate tuples,
where $\Delta_k$, $A_k$, and $R_k$ are as in
Theorem~\ref{thm:finite-observable-image}.  Every coordinate
$p_k=a_k/q_k$ of a rational zero, in reduced form, obeys
\[
 q_k\mid\Delta_k,
 \qquad |a_k|\le R_k,
 \qquad \max(|a_k|,q_k)\le\max(R_k,\Delta_k).
\]
No linear independence of the $b_i$ and no pairwise commutation of
the matrices is required by the checker or the proof.
\end{theorem}

\begin{proof}
At a common zero $p$, substitute its coordinates into every checked
polynomial identity.  The right sides vanish, leaving precisely
\eqref{eq:observable-unit} and \eqref{eq:observable-action} with
$u_k(p)=p_k$.  The finite observable image theorem yields the full
coordinate grid and the reduced-height bounds.

Now enumerate all integer tuples $(z_1,\ldots,z_r)$ with
$-R_k\le z_k\le R_k$, replace them by
$(z_1/\Delta_1,\ldots,z_r/\Delta_r)$ in canonical rational form,
and retain exactly those tuples for which each $F_\ell$ evaluates
to zero.  Every retained tuple is a common zero by the retention test.
Every common zero is visited by the grid bound and retained, proving
completeness.  Distinct integer tuples give distinct rational tuples
because each denominator $\Delta_k$ is fixed and nonzero.  Thus no
additional mathematical selection or deduplication principle is needed.

To implement the Cartesian product with uniform arity, put
$B_k=2R_k+1$ and $B=\prod_kB_k$.  For each $0\le n<B$, recover
its mixed-radix digits $d_k<B_k$ by successive remainder and quotient
operations, and set $z_k=d_k-R_k$.  These arithmetic operations and
the finite decoding loop are PR.  Polynomial evaluation, rational
normalization, list filtering, and the certificate check are also
fixed PR algorithms.  On an invalid certificate the total program
returns a failure flag; its soundness and completeness assertion is
conditional on the checker accepting.

Only the individual eigenvector equations were used.  Matrices that
fail to commute on unused linear directions can therefore still be
valid certificates, and redundant generators cause no difficulty.
\end{proof}

The theorem does \emph{not} find a certificate for an arbitrary input
system.  Certificate production must come from a supplied algebraic
construction, a separately bounded elimination algorithm, or an
independently checked finite output.  A general curve cannot possess
such a certificate for all of its coordinates when it has infinitely
many rational points.  The theorem does not turn a finite rational-point
cardinality assertion into a height bound for that curve.

\paragraph{A stronger algebraic obstruction for a whole plane curve.}
In fact, finite rational-point cardinality need not make the certificate
exist at all.  Let $F\in\mathbb Q[X,Y]$ be monic of positive degree
$m$ in $Y$.  Then $\mathbb Q[X]$ injects into
$\mathbb Q[X,Y]/(F)$.  To prove this without any geometric theorem,
suppose $q(X)=H(X,Y)F(X,Y)$.  If $H\ne0$, its leading
$Y$-coefficient survives multiplication by the monic $F$, so the
right side has $Y$-degree $\deg_YH+m\ge1$; a nonzero $q(X)$
has $Y$-degree zero.  Therefore $H=0$ and $q=0$.

A unit-and-action certificate of size $D$ for the sole equation
$F=0$ would span the entire quotient over $\mathbb Q$: use the
unit vector $c$ for the class of $1$, then the action matrices
to reduce each successive variable multiplication, and add the
reductions for a polynomial's finitely many monomials.  In particular,
$1,X,\ldots,X^D$ would reduce to $D+1$ vectors of length $D$.
Finite rational elimination gives a nontrivial dependence between
them, yielding a nonzero $q(X)\in(F)$ and contradicting the preceding
paragraph.  Thus no such full-coordinate certificate exists, even
when the curve has only finitely many rational points.  For example,
$Y^2+X^2+1=0$ has no rational points, by the order properties of
$\mathbb Q$, yet its monic plane-curve quotient still has this
obstruction.  This polynomial-identity certificate interface handles
finite algebraic sections and systems with supplied finite spanning
data; rational finiteness
of a positive-dimensional algebraic curve is a different kind of
information.  All parts of this obstruction can be read finitely:
the contradiction uses only $D+1$ coefficient vectors, one rational
dependence, and one polynomial degree comparison.

The earlier size-$D$ zero-count theorem remains available.  Its
interpolation argument extends to the relaxed unit condition by using
$c$ as the starting vector for the reduction of the constant polynomial.
Consequently the returned list has at most $D$ entries, although the
coarser grid can be larger.  The grid is a search bound; the algebraic
dimension bound is a stronger count bound.  They serve different roles.

\subsection{A constructor from separated annihilator identities}
\label{sec:separated-annihilators}

The certificate interface becomes an algorithm rather than a mere input
specification whenever one has separated polynomial annihilators with
checked ideal-membership witnesses.  Suppose, for each $1\le k\le r$,
the input supplies a monic univariate polynomial
\[
 q_k(T)=T^{d_k}+\sum_{i=0}^{d_k-1}a_{ki}T^i,
 \qquad d_k\ge1,
\]
and multivariate polynomials $h_{k\ell}$ satisfying
\[
 q_k(X_k)=\sum_{\ell=1}^s h_{k\ell}(X)F_\ell(X).
 \label{eq:separated-annihilator}
\]
These are exact finite polynomial identities.  Merely knowing that
$q_k$ vanishes at some observed solutions is insufficient.

\begin{theorem}[Constructing the spanning certificate]
\label{thm:annihilator-to-module}
From the identities \eqref{eq:separated-annihilator}, one fixed
primitive recursive constructor produces a valid unit-and-action
certificate of size
\[
                         D=\prod_{k=1}^r d_k.
\]
Consequently it produces a complete rational-point list for the
original system.  If $\Delta_k$ clears the denominators of
$a_{k0},\ldots,a_{k,d_k-1}$, a valid coordinate grid uses
\[
 R_k=\max\left(\Delta_k,
               \sum_{i<d_k}|\Delta_k a_{ki}|\right).
 \label{eq:annihilator-grid-radius}
\]
\end{theorem}

\begin{proof}
Index the spanning family by the exponent box
\[
 E=\{e=(e_1,\ldots,e_r):0\le e_k<d_k\},
 \qquad b_e=X_1^{e_1}\cdots X_r^{e_r}.
\]
Use a mixed-radix enumeration of $E$, with the all-zero exponent vector
first.  The unit certificate is $b_0=1$.

Fix $k$ and $e$.  If $e_k<d_k-1$, multiplication by $X_k$
simply sends $b_e$ to the next basis monomial $b_{e+\varepsilon_k}$.
Its matrix column has one entry one and all other entries zero, and
the action residual is identically zero.

If $e_k=d_k-1$, put $\widehat b_e=\prod_{j\ne k}X_j^{e_j}$.
Then
\[
 X_kb_e
   -\left(-\sum_{i<d_k}a_{ki}b_{e[k\leftarrow i]}\right)
   =\widehat b_e q_k(X_k)
   =\sum_{\ell=1}^s
       (\widehat b_e h_{k\ell})F_\ell.
\]
Thus the column of $M_k$ has entry $-a_{ki}$ at the distinct rows
$e[k\leftarrow i]$, with zero elsewhere; the displayed products
$\widehat b_eh_{k\ell}$ are its explicit residual witnesses.
No repeated reduction or termination theorem is needed.  There is
one case distinction and one polynomial multiplication per boundary
column.

After multiplication by $\Delta_k$, a non-boundary column has
absolute sum $\Delta_k$, and a boundary column has absolute sum
$\sum_{i<d_k}|\Delta_ka_{ki}|$.  Hence
\eqref{eq:annihilator-grid-radius} bounds every column sum.
When $d_k=1$ there are no non-boundary columns, so the first term
of the maximum can be omitted; retaining it is harmless.
Apply Theorem~\ref{thm:module-enumeration}.

All exponent vectors, matrix entries, and residual witness polynomials
are produced by finite loops over $E$, the coordinate list, and the
input coefficient lists.  Mixed-radix decoding, multiplication of
polynomials, and the product $D$ are PR uniformly in their finite
input codes.  The correctness invariant is precisely the displayed
action identity for every already generated column.  Its quantifiers
are bounded by the constructed finite dimensions, so Booleanization
gives the required iPRA induction formula.
\end{proof}

The input membership witnesses carry the real elimination information.
The theorem does not assert that such witnesses can be found by an
unbounded search.  It does prove a uniform bounded constructor from
them.  If a supplied identity already has the form
$1=\sum_\ell h_\ell F_\ell$, the correct special case returns the
empty list immediately: substitution at a hypothetical common zero
would give $1=0$.

\paragraph{Optional converse through one adjugate identity.}
The two algebraic interfaces are equivalent as certificate formats,
although the direct coordinate-grid proof needs no determinant.
Indeed, suppose a unit-and-action certificate is given.  Fix $k$,
put $C= X_k I-M_k$, and put $q_k(T)=\det(TI-M_k)$.
The action identities say, for the row $b=(b_1,\ldots,b_D)$,
\[
 bC=r,\qquad r_j=\sum_\ell h_{kj\ell}F_\ell.
\]
Multiply by $\operatorname{adj}C$ and use the finite cofactor
identity $C\operatorname{adj}C=(\det C)I$.  Coordinate $j$ gives
\[
 q_k(X_k)b_j
   =\sum_\ell
       \left(\sum_t h_{kt\ell}(\operatorname{adj}C)_{tj}\right)
       F_\ell.
\]
Multiply by $c_j$, sum over $j$, and use the unit identity
$1-\sum_jc_jb_j=\sum_\ell h_\ell F_\ell$.  The result is an
explicit identity $q_k(X_k)=\sum_\ell u_{k\ell}F_\ell$, with
\[
 u_{k\ell}
   =q_k(X_k)h_\ell+
     \sum_j c_j\sum_t h_{kt\ell}(\operatorname{adj}C)_{tj}.
\]
Here $q_k$ is monic of degree $D$.  Finite permutation sums compute
the determinant and cofactors, so this is also a PR certificate
conversion; the usual cofactor cancellation proves its correctness.
The converse is useful when another application expects separate
elimination polynomials.  It is not needed for denominator bounds,
and it may produce a much larger certificate than the original one.

\subsection{Worked example: a finite section of an elliptic curve}

Consider the two integer polynomials
\[
 F_1=4y^2-4x^3+x-4,
 \qquad F_2=2x^2-x.
\]
The first equation alone is the affine cubic
$y^2=x^3-x/4+1$.  The second cuts out a finite section.  Use the
four generators
\[
             (b_1,b_2,b_3,b_4)=(1,x,y,xy).
\]
This basis and its action witnesses can be constructed automatically
from the two separated identities
\[
 x^2-x/2=F_2/2,
 \qquad y^2-1=\bigl(F_1+(2x+1)F_2\bigr)/4.
\]
Theorem~\ref{thm:annihilator-to-module} uses degrees $(2,2)$
and therefore returns precisely four monomial generators.
The unit identity is $b_1=1$.  In the row-vector convention above,
the multiplication matrices are
\[
 M_x=\begin{pmatrix}
 0&0&0&0\\
 1&1/2&0&0\\
 0&0&0&0\\
 0&0&1&1/2
 \end{pmatrix},
 \qquad
 M_y=\begin{pmatrix}
 0&0&1&0\\
 0&0&0&1\\
 1&0&0&0\\
 0&1&0&0
 \end{pmatrix}.
\]
For completeness, all nonzero residual identities are
\begin{align*}
 x b_2-\tfrac12 b_2&=\tfrac12 F_2,\\
 x b_4-\tfrac12 b_4&=\tfrac12 yF_2,\\
 y b_3-b_1&=\tfrac14 F_1+\tfrac14(2x+1)F_2,\\
 y b_4-b_2&=\tfrac14 xF_1+\tfrac14x(2x+1)F_2.
\end{align*}
The other four action residuals are identically zero.  Expanding the
third identity gives $y^2-1$ on both sides, and multiplying it by
$x$ gives the fourth.  Thus the certificate is a collection of eight
short polynomial identities, not an assumption about a quotient field.

Take $\Delta_x=2$ and $\Delta_y=1$.  The absolute column sums of
$A_x=2M_x$ are $(2,1,2,1)$, and those of $A_y=M_y$ are all one.
Hence $R_x=2$ and $R_y=1$.  Every rational common zero lies among
the fifteen candidate points
\[
 x\in\{-1,-1/2,0,1/2,1\},
 \qquad y\in\{-1,0,1\}.
\]
Filtering the two equations gives exactly
\[
                 (0,1),\ (0,-1),\ (1/2,1),\ (1/2,-1).
\]
One can verify the last filter symbolically as well: $F_2=0$
implies $x=0$ or $x=1/2$ by the integral-domain law; at either
value $F_1=0$ reduces to $y^2=1$, and therefore $y=1$ or $y=-1$.
The point count four attains the size of the spanning family.  The
denominator information is genuinely part of the result: the argument
forces $x$ to have reduced denominator dividing two and forces $y$
to be an integer before any solution is selected.

\subsection{Principal-open conditions and finite chart covers}

A finite algebra certificate can also treat equations together with a
polynomial inequation $G(x)\ne0$.  Introduce one extra rational
coordinate $t$ and the equation
\[
                           tG(x)-1=0.
\]
Every solution of the original system with $G(x)\ne0$ lifts uniquely
to the extended system by $t=1/G(x)$.  Conversely, an extended solution
forces $G(x)\ne0$ and projects to an original solution.  A valid
unit-and-action certificate for the extended system therefore gives
a complete finite list for that principal-open part.  Projection may
be followed by a bounded exact deduplication pass, although uniqueness
of $t$ already prevents repetitions coming from a single chart.

For projective equations, use the disjoint charts in which coordinate
$j$ is the first nonzero coordinate.  Normalize it to one, set all
earlier coordinates to zero, and substitute into the homogeneous
equations.  A supplied finite algebra certificate for every resulting
affine system gives a complete projective rational-point list: each
projective point belongs to exactly one such chart and has exactly one
normalized representative there.  Thus several small algebraic
certificates can replace one large elimination, without an appeal to
compactness.  This construction is conditional on the actual chart
certificates; no claim of general rational-point decidability is hidden
in the chart decomposition.

\subsection{The exact iPRA and Lean obligations}

The complete human proof reduces to finite arithmetic.  A practical
formalization should use the following dependency order.
\begin{enumerate}
\item Define signed integer and reduced rational codes, finite vectors,
matrix multiplication, mixed-radix tuple decoding, and sparse polynomial
evaluation.  Invalid codes receive specified default values.  Prove
normalization and Euclidean arithmetic by quantifier-free induction.
\item Prove a signed finite-vector Bezout lemma: if the gcd of the
absolute coordinates is one, the extended-gcd fold returns a coefficient
vector $e$ with $e\cdot w=1$.  Track the current gcd and the coefficient
identity as a bounded invariant.
\item Prove primitive scaling, then the integer-eigenvalue lemma using
the displayed coordinate divisibility identities.  The bound uses a
finite maximum, the triangle inequality for integer sums, and cancellation
of a positive integer.  No determinant library is required.
\item Define the certificate checker by coefficient equality for
\eqref{eq:module-unit-certificate} and
\eqref{eq:module-action-certificate}.  Prove its substitution lemma:
acceptance plus vanishing of the input equations implies the evaluated
unit and action identities.
\item Prove grid completeness and the finite filter specification.
The grid invariant states that each processed mixed-radix index has
been decoded correctly and retained exactly when all equation tests
succeed.  The input rational point is a free parameter when proving
that its grid index is visited.
\item Give explicit PR syntax for these fixed algorithms, compile the
bounded invariants to quantifier-free formulas, and build actual
\texttt{Prf PRAAx} derivations.  Only then apply the repository's
existing translation theorem into the native scalar presentation.
\end{enumerate}

The accompanying \path{lean/FiniteAlgebraTargets.lean} supplies
fully defined sparse polynomial operations, the certificate checker,
the Cartesian-grid producer, and the exact enumeration target.
Section~\ref{sec:formalization} gives the complete dependency order,
distinguishing checked host lemmas from the final object derivation.

The general observable theorem is a schema for fixed PR observables
with supplied local proofs.  The polynomial specialization is uniform
over arbitrary finite polynomial and certificate codes because its
evaluator executes one fixed polynomial grammar.  These two kinds of
uniformity must not be confused with a total universal evaluator for
all primitive recursive functions.

\section{From a nonzero algebraic norm to a complete rational intersection list}
\label{sec:norm-enumeration}

The earlier norm packets give sharp bounds on the number of distinct
points in a supplied list \cite{Demand,Coherent}. For an \emph{explicitly supplied section},
their finite data yield more: a complete rational-point list, with a
primitive recursive coordinate bound.  The additional ingredient is
the elementary rational-root argument.  It is useful to state this
separately, because a list of actual sections and the assertion that
every tuple admits some sections have different computational content.

\subsection{A total bounded rational-root procedure}

Encode a rational as a signed numerator and a positive denominator in
lowest terms, and a rational polynomial as a finite coefficient list.
The zero polynomial is detected by a bounded coefficient scan.
For a nonzero polynomial \(H\in\mathbb Q[T]\), multiply by the positive
product of all coefficient denominators.  This gives a nonzero integer
polynomial \(h\), with precisely the same rational zeros.  Scan for its
first and last nonzero coefficients and write
\[
 h(T)=T^e(c_0+c_1T+\cdots+c_dT^d),
 \qquad c_0c_d\ne0.
\]
Here \(d\) may be zero.  Put
\[
 B(H)=\max\{1,|c_0|,|c_d|\}.
 \tag{RR1}\label{eq:norm-root-bound}
\]
Define \(\operatorname{Roots}_{\mathbb Q}(H)\) by scanning all pairs
\[
 -B(H)\le p\le B(H),\qquad 1\le q\le B(H),\qquad
 \gcd(|p|,q)=1,
\]
and retaining \(p/q\) exactly when \(H(p/q)=0\).  Each retained
rational has a unique representation in this scan.  The totalized
procedure returns a failure tag on the zero polynomial; it does not
return an allegedly complete finite list in that case.

\begin{lemma}[Bounded rational roots]\label{lem:norm-rational-roots}
For every nonzero rational polynomial \(H\), the displayed procedure
returns exactly its rational roots, without repetition.  Every root
has reduced numerator and denominator bounded in absolute value by
\(B(H)\).  The procedure and its correctness proof use only primitive
recursion and quantifier-free induction on finite arithmetic data.
\end{lemma}

\begin{proof}
Soundness follows from the final exact substitution test.  Let
\(p/q\), with \(q>0\) and \(\gcd(|p|,q)=1\), be a root.
If \(p=0\), reducedness gives \(q=1\), so the scan includes it.
Suppose \(p\ne0\).  Cancelling the nonzero factor \((p/q)^e\)
and multiplying by \(q^d\) gives
\[
 c_0q^d+c_1pq^{d-1}+\cdots+c_dp^d=0.
 \tag{RR2}\label{eq:norm-cleared-root}
\]
For \(d=0\), this contradicts \(c_0\ne0\), so there is no such
root.  For \(d>0\), reduction modulo \(p\) shows
\(p\mid c_0q^d\), and reduction modulo \(q\) shows
\(q\mid c_dp^d\).  Euclid's lemma, applied to the coprime pairs
\((p,q^d)\) and \((q,p^d)\), gives
\(p\mid c_0\) and \(q\mid c_d\).  The coefficients are nonzero,
so \(|p|\le|c_0|\) and \(q\le|c_d|\).  Thus every root occurs in
the scan.

For an elementary arithmetic implementation of the cancellation step,
Bezout coefficients for \(p,q\) are computed by a Euclidean algorithm
with fuel bounded by \(|p|+q+1\).  Multiplying the Bezout relation
finitely many times gives the required coprimality of powers.
All integer divisions, signs and zero tests are decidable.  There are
at most \((2B(H)+1)B(H)\) candidate pairs.  Degree, the coefficient
identities, and the scan prefix are finite arrays; correctness is
therefore proved by Booleanized bounded invariants.  No assertion
about an infinite set of roots is used as an induction formula.
\end{proof}

The bound is deliberately coarse.  Divisor enumeration, content
removal and synthetic division can substantially reduce the computation,
but are unnecessary for the primitive recursive theorem.  A nonzero
constant polynomial gives the empty list.  A polynomial \(cT^e\),
with \(e>0\), gives exactly \(0\).  Thus neither zero constant terms
nor constant quotients require an exception left outside the algorithm.

\subsection{The norm and the intersection algorithm}

Let
\[
 F(X,Y)=Y^m+f_{m-1}(X)Y^{m-1}+\cdots+f_0(X),
 \qquad m\ge1,
\]
with \(f_j\in\mathbb Q[X]\), and let \(G\in\mathbb Q[X,Y]\)
be arbitrary.  Monic division in \(Y\) identifies the quotient by
\(F\) with a free \(\mathbb Q[X]\)-module generated by
\(1,Y,\ldots,Y^{m-1}\).  Reduce each \(GY^j\) modulo \(F\),
and let \(M_G(X)\) be the matrix whose \(j\)-th column is that
coefficient vector.  Define the polynomial norm
\[
 N_G(X)=\det M_G(X)\in\mathbb Q[X].
\]
Reduction and the determinant can be computed by finite coefficient
operations.  In particular, scanning the \(m^m\) length-\(m\)
index tuples and retaining the permutations is an explicit primitive
recursive determinant implementation.  No general algorithm for ideal
membership, algebraic closure or factorization is needed.

\begin{lemma}[Common zeros give norm zeros]
\label{lem:norm-intersection-vanishing}
If \((a,b)\in\mathbb Q^2\) satisfies \(F(a,b)=G(a,b)=0\), then
\(N_G(a)=0\).
\end{lemma}

\begin{proof}
Put \(v=(1,b,\ldots,b^{m-1})\).  Evaluating each monic reduction
identity at \((a,b)\) gives
\[
 vM_G(a)=G(a,b)v=0.
\]
The adjugate identity, proved by finite determinant expansion, yields
\(vM_G(a)\operatorname{adj}(M_G(a))=N_G(a)v\).  Its first
coordinate is \(N_G(a)\), because \(v_0=1\).  It therefore vanishes.
This proof allows repeated roots, singular fibers and \(m=1\).
\end{proof}

Assume now that the finite coefficient test verifies \(N_G\ne0\).
First compute the finite list
\[
 A=\operatorname{Roots}_{\mathbb Q}(N_G).
\]
For each \(a\in A\), form \(F_a(Y)=F(a,Y)\), which is monic
of degree \(m\), hence nonzero.  Compute
\(B_a=\operatorname{Roots}_{\mathbb Q}(F_a)\).  The output list is
\[
 \mathcal I(F,G)
 =\bigl[(a,b):a\in A,\ b\in B_a,\ G(a,b)=0\bigr].
 \tag{NI}\label{eq:norm-intersection-list}
\]
All tests are exact rational tests.  A norm root can be extraneous;
the last test is essential.  The empty list is a permitted result.

\begin{theorem}[Complete rational intersection enumeration]
\label{thm:norm-complete-intersection}
For every monic \(F\) as above and every \(G\) with \(N_G\ne0\),
\(\mathcal I(F,G)\) is a duplicate-free complete list of the
common rational zeros of \(F,G\).  It is produced by one primitive
recursive algorithm, uniformly in the finite coefficient arrays and
the numerical degree \(m\).

An explicit primitive recursive height bound is
\[
 H(F,G)=\max\Bigl(\{1,B(N_G)\}
       \cup\{B(F_a):a\in A\}\Bigr).
 \tag{NH}\label{eq:norm-intersection-height}
\]
Here the height of a rational pair is the maximum of the absolute
values of its reduced numerators and its positive denominators.
Every common rational zero has height at most \(H(F,G)\).
\end{theorem}

\begin{proof}
If a pair is output, its first coordinate is in \(A\), its second
coordinate is an exact root of \(F_a\), and the final test gives
\(G(a,b)=0\).  Thus it is a common zero.  Conversely, any common
zero \((a,b)\) satisfies \(N_G(a)=0\) by
Lemma~\ref{lem:norm-intersection-vanishing}; completeness of the first
root scan puts \(a\) in \(A\).  Since \(F_a(b)=0\), completeness
of the second scan puts \(b\) in \(B_a\).  The final test passes.
Each root scan uses canonical fractions without repetitions, and
the first coordinates distinguish the blocks, so the pair list has
no repetitions.  The height assertion follows directly from
Lemma~\ref{lem:norm-rational-roots} in the two scans.

The number of first-coordinate candidates is bounded by
\((2B(N_G)+1)B(N_G)\).  After this finite list is computed, the
maximum of its specialized bounds is another primitive recursive
fold.  These numerical bounds govern the second scans.  Computing a
bound from a previously computed finite list is ordinary composition
and primitive recursion; it is not unbounded minimization.  Polynomial
evaluation and rational normalization are fixed primitive recursive
programs on coefficient and point codes.  Soundness and completeness
are obtained from bounded scan invariants and the displayed finite
identities.  Thus the written argument supplies an iPRA proof strategy
without induction over an unbounded root predicate.
\end{proof}

The requirement \(N_G\ne0\) is a sufficient certificate, not a
necessary characterization of finite rational intersections.
For example, a zero norm can arise from a common algebraic component
with no rational points.  The algorithm reports an unsupported input
when the norm is identically zero; it draws no conclusion that the
rational intersection is infinite.  In particular \(G=0\) is rejected.
No irreducibility or nonsingularity hypothesis is needed on supported
inputs.  The integer points are obtained simply by retaining output
coordinates whose denominators are one.

\subsection{The norm also constructs a finite multiplication certificate}
\label{sec:norm-to-module}

The norm interface can be connected directly to the finite-module
enumeration theorem.  It supplies an explicit certificate constructor
for the entire supported monic-plane class, not merely an independent
root-finding algorithm.

\begin{proposition}[A norm-to-module certificate constructor]
\label{prop:norm-to-module}
Given \(F,G\) as above with nonzero norm of degree \(d>0\), a
primitive recursive procedure constructs a unit-and-action certificate
for the ideal \((F,G)\), using the \(md\) monomials
\[
             X^iY^j\qquad(0\le i<d,\ 0\le j<m).
\]
If \(N_G\) is a nonzero constant, it instead constructs an ideal
identity for \(1\), certifying that there are no common zeros.
\end{proposition}

\begin{proof}
Retain the quotients from the monic reductions used to construct
\(M_G\): for \(0\le j<m\),
\[
 GY^j=\sum_{i<m}(M_G)_{ij}Y^i+Fh_j.
\]
For the rows \(b=(1,Y,\ldots,Y^{m-1})\) and
\(h=(h_0,\ldots,h_{m-1})\), this says
\(bM_G=Gb-Fh\).  Right multiplication by the adjugate gives
\[
 N_G b=G b\operatorname{adj}(M_G)
           -Fh\operatorname{adj}(M_G).
\]
Its first coordinate is an explicit polynomial identity
\[
 N_G(X)=U(X,Y)G(X,Y)+V(X,Y)F(X,Y),
 \tag{NM}\label{eq:norm-membership}
\]
where \(U=(b\operatorname{adj}(M_G))_0\) and
\(V=-(h\operatorname{adj}(M_G))_0\).  Thus no ideal-membership
oracle is needed.  If \(N_G=c\ne0\) is constant, division of
\eqref{eq:norm-membership} by \(c\) gives the claimed identity
for \(1\).

Otherwise divide \(N_G\) by its nonzero rational leading
coefficient to obtain a monic \(q(X)\) of degree \(d\), retaining
the correspondingly scaled witnesses in \eqref{eq:norm-membership}.
Use the displayed monomial rectangle as the generating family, with
\(1\) its first member.  Multiplication by \(X\) stays inside
the rectangle unless \(i=d-1\); in that case the monic relation
\(q(X)\), multiplied by \(Y^j\), gives the action coefficients
and their ideal witnesses directly.

Multiplication by \(Y\) stays inside the rectangle unless
\(j=m-1\).  In that case write
\[
 -X^i f_t(X)=s_t(X)q(X)+r_t(X),\qquad \deg r_t<d
 \quad(0\le t<m)
\]
by bounded monic division in \(X\).  The exact identity
\[
 YX^iY^{m-1}-\sum_{t<m}r_t(X)Y^t
       =X^iF+\sum_{t<m}s_t(X)q(X)Y^t
\]
provides this boundary column.  Substitute the already retained ideal
witness for \(q\) on the right.  All remaining terms are explicit
multiples of \(F,G\), and the left coefficient combination lies
inside the rectangle.  These are precisely the required unit and
action identities.

The division in \(Y\) has its input-degree bound, the determinant
and adjugate are finite permutation sums, and each subsequent monic
division in \(X\) has its input-degree bound.  The constructor
performs these fixed loops over \(md\) monomials and retains the
identities at each step.  Its proof uses coefficientwise finite
identities as bounded invariants.  It is therefore primitive recursive
uniformly in the supplied polynomial codes and has the indicated
iPRA correctness proof strategy.
\end{proof}

The rational-root scan is usually smaller than the multiplication
certificate produced by this constructor.  Both have value: the scan
gives a direct implementation, while the constructor lets the same
general proof-producing finite-algebra interface certify the result.

\subsection{Existing weighted norms become effective section lists}

The previously supplied norm theorems treat, among other cases,
\[
 F(X,Y)=Y^m-cX^d+\text{terms of smaller weight},
 \quad \mathrm{wt}(X)=m,\quad\mathrm{wt}(Y)=d,
\]
where \(m\ge2\), \(d\ge1\), \(\gcd(m,d)=1\) and \(c\ne0\).
For a nonzero reduced polynomial \(G\), with \(\deg_YG<m\), those
theorems give an explicit nonzero leading coefficient for \(N_G\).
The new enumeration theorem uses precisely this nonvanishing output;
it need not reprove the sharper norm degree or fiber multiplicity bound.

Suppose a proper rational linear subspace in a finite monomial
evaluation space is supplied by a finite matrix.  Bounded rational
elimination produces a nonzero annihilator \(\lambda\); the
corresponding polynomial \(G_\lambda\) is a nonzero linear
combination of distinct reduced monomials.  The weighted theorem
certifies \(N_{G_\lambda}\ne0\), so
Theorem~\ref{thm:norm-complete-intersection} lists every rational
curve point on that section.  If only one annihilator is used, the
list may include points in the containing hyperplane but outside
the original subspace.  A final exact subspace-membership test
removes them.  Hence the section list is exact, not merely a cover
by candidates.

As a small example, take
\[
 F=Y^2-X^3+X-1,\qquad G=X^2-X.
\]
Since \(G\) has no \(Y\)-dependence, \(M_G=G I_2\), whence
\(N_G=(X^2-X)^2\).  Its rational-root list is \((0,1)\).
At either first coordinate, \(F_a(Y)=Y^2-1\), with roots
\(-1,1\).  The last filter retains exactly
\[
 (0,-1),\ (0,1),\ (1,-1),\ (1,1).
\]
This is the sharp four-point example already used for counting in
the older framework; here the same finite section data give a
certified exhaustive search.  The argument establishes a new
connection between existing packet interfaces, not priority for
the underlying rational-root or elimination methods.

\subsection{Global covers, tuple-local covers, and bounded coefficients}

Let \(S(a,b)\) be any fixed decidable condition on rational points
of \(F=0\), such as integrality.  Suppose actual polynomials
\(G_1,\ldots,G_s\) are supplied, each with nonzero norm, together
with a proof of the single global assertion
\[
 \forall a,b\quad
 F(a,b)=0\ \wedge\ S(a,b)
 \ \Longrightarrow\ \bigvee_{j=1}^{s}G_j(a,b)=0.
 \tag{GC}\label{eq:norm-global-cover}
\]
The cover assertion can be a supplied object-theory derivation;
merely checking its truth on an initial search box is insufficient.

\begin{corollary}[A supplied global cover gives all points]
\label{cor:norm-global-list}
Under \eqref{eq:norm-global-cover}, concatenate the lists
\(\mathcal I(F,G_j)\), remove duplicate pairs and retain the
pairs satisfying \(S\).  The result is a complete finite list of
all rational points satisfying \(F=0\) and \(S\).
The maximum of the bounds \(H(F,G_j)\), or one for an empty
cover, is a height bound.
\end{corollary}

\begin{proof}
Every retained pair is a curve point satisfying \(S\).  Conversely,
a point satisfying these conditions belongs to one supplied section
by the finite disjunction in \eqref{eq:norm-global-cover}.  The
corresponding intersection theorem puts it in that section's list.
Concatenation and deduplication preserve membership.  This is a
finite intuitionistic disjunction elimination, not a selection of
an infinite family of exceptional subspaces.  If the cover is
empty, its global assertion already proves that no such point exists.
\end{proof}

The Siegel route in the repository instead asks that \emph{each
supplied finite tuple} admit some cover by at most \(M\) proper
subspaces.  Its sections may change with the tuple, and their
coefficients may depend on the tuple's coordinates.  Applying the
present theorem to a cover of an already supplied tuple locates
only points of that cover; it gives no guarantee that all further
points lie there.  Even a sharp bound on the number of points in
each section does not repair this quantifier change.  Thus the
new enumeration theorem resolves the localization of explicit
sections, but does not supply the missing global approximation
cover or prove general Siegel or Faltings.

A useful sufficient strengthening can nevertheless be stated
precisely.  Fix a list of \(r\) distinct reduced monomials and a
positive integer \(C\).  There are at most
\((2C+1)^r-1\) nonzero integer coefficient vectors
\(\lambda\) with \(|\lambda_i|\le C\).  For each vector, compute
the norm of its monomial combination and retain those with nonzero
norm.  If a proof establishes that every point in \(S\) lies on
one of these retained sections, the preceding corollary computes
all such points.  In the weighted coprime setting, every nonzero
combination is automatically retained.  Consequently a
\emph{uniform coefficient-height bound} for the exceptional
annihilators would suffice, even if no preferred global cover were
selected.  Tuple-dependent coefficient bounds do not suffice.
This identifies a concrete missing quantitative input rather than
an unspecified need for choice.

For Mordell--Weil, none of these conclusions constructs a rational
point on every relevant descent covering.  Such covering curves are
typically positive-dimensional, and a finite quotient or nonzero
intersection norm becomes available only after an additional
independent equation is supplied.  Solving its resulting finite
intersection does not prove that some intersection contains a
representative of every required quotient class.  The rational
representative-list obligation therefore remains intact.

\subsection{Lean and object-iPRA obligations}

The definitions needed for a complete formalization are finite:
signed rational codes with normalization; univariate and bivariate
coefficient arrays; monic \(Y\)-reduction; a determinant fold;
\(\operatorname{RootBound}\), \(\operatorname{RootList}\), and
\(\operatorname{IntersectionList}\); and exact Boolean root and
subspace checks.  The recommended lemma order is:
\begin{enumerate}
\item prove gcd/Bezout cancellation, the denominator-cleared
polynomial identity, and rational-root bounds;
\item prove the bounded root scan's soundness and completeness;
\item prove monic reduction under evaluation and the adjugate
identity, then the norm-vanishing lemma;
\item combine the two root scans into complete intersection
enumeration, with an explicit maximum over first-coordinate candidates;
\item combine finitely many intersection lists under a supplied
derivation of the global cover formula.
\end{enumerate}
At the embedded level, each scan's invariant states that its output
contains exactly the successful candidates examined so far, as a
bounded Boolean equality.  Coefficient clearing, reduction and
determinant identities are equalities of finite codes or bounded
coefficient arrays.  Euclidean descent uses fixed numerical fuel.
Free point codes remain parameters throughout these inductions.
Only after constructing the arithmetic derivations does one use
the repository's native equivalence.  An ambient Lean theorem
about these arrays, or a type-checked declaration of the target,
does not on its own construct the required \texttt{Prf PRAAx} term.

\section{How the framework composes across mathematical domains}
\label{sec:applications}

The simulation component and the finite-algebra component have different
jobs. Simulation preserves a proved observation during a finite
construction. Finite observable modules restrict the possible rational
values of an observation. Reconstruction turns the resulting finite
possibilities into actual objects when a finite fiber procedure is
available. The following examples show how to use the three steps
without adding an implicit existence principle.

\subsection{Modular computation and exact polynomial evaluation}

Fix $q\ge1$ and let $\rho_q(a)$ be the remainder of $a$ on division
by $q$. Define the source step $f(a)=a^2$ and the target step
$g(r)=\rho_q(r^2)$. The local fusion law is
\[
                  \rho_q(a^2)=\rho_q(\rho_q(a)^2).
\]
To prove it arithmetically, write $a=qA+r$, $0\le r<q$.
Expanding gives $a^2=q(qA^2+2Ar)+r^2$. Division of $r^2$ by
$q$ therefore gives the same final remainder as division of $a^2$;
uniqueness follows by subtracting two division equations with remainders
between zero and $q-1$. These are finite integer equations and
inequalities, with a quantifier-free induction proof for division.
Theorem~\ref{thm:fusion} now transports the entire repeated-squaring
construction through $\rho_q$.

The statement is uniform in the numerical modulus. To instantiate the
unary theorem literally, use state $(q,a)$ and a representation map
$(q,a)\mapsto(q,\rho_q(a))$. The guard $q\ge1$ belongs in the
relational invariant; a total default for $q=0$ is never used to
justify modular mathematics. This pattern applies equally to computations
in a presented number ring after its finite arithmetic and reduction
laws have been supplied.

For a rational polynomial $P$, evaluation at a fixed rational tuple
$p$ preserves addition and multiplication. The multiplication assertion
has a direct finite proof: distribute the two coefficient sums, multiply
the monomials, and use $p^{\alpha+\beta}=p^\alpha p^\beta$.
There are only the finitely many coefficient pairs in the input lists.
The local proof transports repeated squaring and more general finite
expression evaluation into rational arithmetic. It can also transport
coefficient identities into equations at a common zero, the precise
bridge used by the certificate theorem.

\subsection{Truncated series and coherent exact computation}

Let $\pi_N$ discard all coefficients above degree $N$. For finite
polynomials or coefficient lists, convolution gives
\[
      \pi_N(AB)=\pi_N\bigl(\pi_N(A)\pi_N(B)\bigr).
\]
Indeed the coefficient of degree $k\le N$ is
$\sum_{i=0}^{k}a_i b_{k-i}$. Every contributing index is at most
$N$, so truncating either input leaves it unchanged. For $k>N$
both truncated outputs have zero coefficient. Coefficient equality
is a bounded list condition and hence an eligible arithmetic invariant.

Iteration fusion applies to repeated polynomial operations whose local
truncation law has this form. It yields the exact requested jet without
requiring a completed infinite series as an object of \ipra. This
connects to the repository's filtered-lifting and formal-series work
\cite{Coherent}. Differentiation uses a different demand: coefficients
through $N+1$ determine the derivative through $N$. A correct
simulation must include that shift rather than claiming that arbitrary
operations preserve the same truncation level.

Exact coherence does not imply an analytic tail bound. A convergent
analytic interpretation must supply a PR function $T(N,k)$ or another
explicit estimate making the omitted tail small enough at requested
precision $k$. The present transport theorem then handles the finite
coefficient calculation; the tail theorem remains its own mathematical
input. This is a common interface for the Gamma, zeta, calculus,
Lagrange inversion and regular local-branch developments, rather than
a replacement for their analytic estimates.

\subsection{Redundant algebraic presentations and observational equality}

The point-enumeration checker intentionally does not require a basis or
commuting matrices. This is useful, not just a technical weakening.
Consider
\[
 F_1=x^2-x,\qquad F_2=2y-x,\qquad (b_1,b_2,b_3)=(1,x,x).
\]
Use
\[
 M_x=\begin{pmatrix}0&0&0\\1&1&1\\0&0&0\end{pmatrix},
 \qquad
 M_y=\begin{pmatrix}0&0&0\\0&0&0\\1/2&1/2&1/2\end{pmatrix}.
\]
They do not commute: the nonzero row of $M_xM_y$ is its second,
whereas that of $M_yM_x$ is its third. Yet the action identities
are valid. For the $x$ action the residuals are $0,F_1,F_1$.
For the $y$ action they are
\[
 \tfrac12F_2,\qquad
 \tfrac12F_1+\tfrac x2F_2,\qquad
 \tfrac12F_1+\tfrac x2F_2.
\]
The unit generator is $b_1=1$. Hence this is an accepted certificate.
It gives $\Delta=(1,2)$ and $R=(1,1)$, and the grid filter returns
exactly $(0,0)$ and $(1,1/2)$.

The apparent matrix discrepancy acts in a redundant direction. At any
common zero, vectors of coefficients $c,d$ representing the same
value are related by
\[
                  \sum_i c_i b_i(p)=\sum_i d_i b_i(p).
\]
That relation, with the common-zero guard and the fixed point $p$
stored as parameters, is the right observation to preserve. Equality
of coefficient vectors would be unnecessarily strong. A relational
simulation therefore allows a convenient redundant presentation without
first constructing a canonical basis. This matters in weak arithmetic,
where basis selection, rank certificates and coefficient growth are
additional obligations which the desired conclusion may not need.

\subsection{A finite open part of a positive-dimensional system}

Equations plus guards occur in rational maps, affine charts, and
elimination with nonzero denominators. Here is a complete concrete
example of the localization principle.
The equations
\[
                xy=0,\qquad x(x-1)=0
\]
have infinitely many rational solutions: the entire line $x=0$ is
included. On the open part $x\ne0$, the only solution is $(1,0)$.
Introduce $t$ and use
\[
              F_1=xy,\quad F_2=x(x-1),\quad F_3=tx-1.
\]
The one-element spanning family $(1)$ has action matrices $M_x=(1)$,
$M_y=(0)$, $M_t=(1)$. The complete certificate is
\begin{align*}
 x-1&=tF_2-(x-1)F_3,\\
 y&=tF_1-yF_3,\\
 t-1&=-t^2F_2+\bigl(1+t(x-1)\bigr)F_3.
\end{align*}
Each identity follows by expanding both sides; together with the unit
identity they are all the checker needs. The coordinate radii are
$(1,0,1)$, so nine candidate triples suffice. Filtering gives precisely
$(1,0,1)$; projection gives the desired open-part list.

This example illustrates a concrete certificate producer in a useful
special case. It does not infer that localization always makes a system
finite. The guard is explicit, its inverse coordinate is unique, and
the extended system has its own verified action data. Several guards
$G_i\ne0$ can be combined as $G=\prod_iG_i\ne0$ over $\Q$,
or treated with separate inverse coordinates.

\subsection{Why finite rational sets on curves are a different problem}

A multiplication certificate is stronger than the assertion that the
rational zero set is finite. The certificate identities imply that the
entire quotient algebra $\Q[X_1,\ldots,X_r]/(F_1,\ldots,F_s)$
is spanned by finitely many classes: start with the unit combination,
multiply by variables using the action tables, and extend by linearity.
This argument applies to all polynomials, without testing rational points.

For a single monic plane equation $F(X,Y)$ of positive degree in $Y$,
the map $\Q[X]\to\Q[X,Y]/(F)$ is injective. A nonzero multiple
$HF$ has $Y$-degree $\deg_Y H+\deg_Y F>0$ and so cannot be a
nonzero polynomial in $X$ alone. Consequently the powers of $X$
remain linearly independent in the quotient. No finite spanning
certificate for both coordinates can exist for that full ideal.
This remains true if the curve has only finitely many rational points.

Thus the present theorem does not furnish a direct decision procedure
for rational points on every curve. It applies to finite algebraic
sections, suitably finite open charts, or other situations for which
the extra finite action data are actually available. The norm constructor
explains exactly how an additional section equation supplies those data.
The missing global arithmetic work is to show that the points of
interest are covered by finitely many such explicit sections, with
uniformly available coefficients or a suitable bound on them.

\subsection{A precise task decomposition for the difficult reports}

\begin{longtable}{P{.20\linewidth}P{.38\linewidth}P{.32\linewidth}}
\toprule
Area & Reusable interface supplied here & Further theorem still required\\
\midrule\endhead
Siegel-style sections & A nonzero norm produces all rational points
in a supplied monic-plane section, and a finite algebra certificate. &
A global cover or bounded coefficients for its sections; the quantitative
approximation theorem is not supplied.\\[4pt]
Faltings-style finiteness & Complete lists for individually certified
finite algebraic pieces; finite union and chart transport. & A method
covering the rational points of a positive-dimensional curve by such
explicit pieces.\\[4pt]
Mordell--Weil descent & Proof transport for finite descent computations
and exact enumeration of separately certified finite loci. & Actual
covering representatives and a proof that they cover the required group
quotient.\\[4pt]
Formal algebra and local expansions & Checked simulation through
truncation and evaluation, once their local object laws are supplied. &
Application-specific local laws, singular branch analysis, and any
analytic tail bounds.\\[4pt]
Analytic number theory & Composable finite exact arithmetic and
explicit certificate checking. & The uniform analytic estimates named
in the prime-distribution report.\\
\bottomrule
\end{longtable}

This division is productive because it names executable and provable
subtasks. A future global-cover theorem can call the enumeration module
directly. A future formal coefficient library can reuse the checked
simulation builder. Until those inputs exist, the completed sublemmas
remain useful without being misreported as proofs of the original
deep theorems.

\section{Lean formalization: checked core and a complete path forward}
\label{sec:formalization}

\subsection{Pinned object syntax and standalone build}

The archive builds with Lean 4.32.0 and its standard library. It includes
unchanged copies of the pinned repository's \code{Syntax.lean},
\code{Deduction.lean}, \code{PR.lean} and
\code{NativeBiInterpretation.lean}. The first three module bodies are
concatenated into \code{ObjectKernel.lean}; their imports are replaced
by \code{import Std}, and a compatibility notation spells the
natural-number type. No object definition, axiom constructor, inference
rule or upstream theorem body is changed. A reproducible provenance
checker verifies the embedded bodies and the pinned Git blob hashes.
The full native bi-interpretation file is retained as a reference and
is not rebuilt by the standalone core.

This arrangement avoids downloading Mathlib merely to use the small
deductive kernel. Integration into ProveIt replaces the compatibility
import by the repository's \code{ConciseIPRA.PR} import and retains
the new namespaces and proofs. The validation record is a standalone
build against the verified upstream bodies, not a claim that the whole
ProveIt workspace has been rebuilt.

\begin{longtable}{P{.29\linewidth}P{.40\linewidth}P{.21\linewidth}}
\toprule
Module & Principal content & Status\\
\midrule\endhead
\path{ObjectKernel.lean} & The three upstream syntax, deduction and
arithmetic module bodies. & Compiled; provenance checked.\\[3pt]
\path{ObjectFusion.lean} & PR iteration equations and the fusion
derivation builder. & Compiled object theorems.\\[3pt]
\path{ObjectSimulation.lean} & Substitution preserves QF status;
relational simulation builder. & Compiled object theorem.\\[3pt]
\path{SquaringTransport.lean} & Repeated-squaring transport and an
unconditional successor instance. & Compiled object theorems.\\[3pt]
\path{PrimitiveEigenvalue.lean} & Finite weighted eigen-equations
and the primitive-vector denominator lemma. & Compiled host integer
theorems.\\[3pt]
\path{EigenvalueBound.lean} & Finite absolute-column estimate and
the eigenvalue bound. & Compiled host integer theorems.\\[3pt]
\path{FiniteAlgebraTargets.lean} & Complete finite programs,
remaining theorem specifications and the filter reduction. & Definitions
compile; main completeness target open.\\[3pt]
\path{FiniteAlgebraExamples.lean} & Exact execution of the elliptic
section certificate and grid. & Compiled diagnostics.\\
\bottomrule
\end{longtable}

There are no custom axioms, opaque placeholders, \code{sorry} or
\code{admit} declarations in the new modules. The target file does
not smuggle the desired theorem into a structure field: the remaining
claims are explicitly named propositions and are not asserted.

\subsection{The completed syntactic theorems}

The important distinction is the type of a theorem's conclusion. In
namespace \path{LeanProofs.ConciseIPRA.FiniteMethods}, the simplified
shape of the checked fusion theorem is
\begin{verbatim}
theorem fusion (f g h : PR 1)
    (hlocal : ([] : List (Fm PTm 1)) |-[PRAAx]
      commutes f g h) :
    ([] : List (Fm PTm 1)) |-[PRAAx]
      Fm.all (fusionInvariant f g h)
\end{verbatim}
Here the printed ASCII turnstile is explanatory; the source uses the
repository's scoped judgment notation. \code{commutes} is precisely
$h(f(x))=g(h(x))$, and \code{fusionInvariant} is the equality in
Theorem~\ref{thm:fusion}. No extra assumption about an evaluation
function is present.

For simulation the inputs are $R:\code{Fm PTm 2}$, a proof of
\code{R.QF}, $f,g:\code{PR 1}$, and the local judgment
$[R]\vdash_{\code{PRAAx}}R(f(x),g(y))$. The output is
$[R]\vdash_{\code{PRAAx}}\forall n\,R(\iter_f(n,x),\iter_g(n,y))$.
The code verifies every lifting of a substitution under a binder.
Both modules use \code{PRA.qfInd}, \code{Prf.subst},
\code{Prf.eqSubst} and the defining equations, exactly as the human
proofs specify.

The type \code{Prf} lives in Lean's proposition universe; this
development proves derivability and supplies a checked derivation term.
It does not serialize a numerical proof trace for the tiny Python
checkers in the other project. Such serialization, if desired, is
another implementation task. This distinction has no effect on the
object judgment proved here.

\subsection{The integer arithmetic already checked}

\code{PrimitiveEigenvalue.lean} stores a finite list of triples
$(e_i,w_i,z_i)$, representing a Bezout coefficient, a primitive-vector
coordinate, and its integral matrix image. It defines
\[
          W=\sum_i e_iw_i,\qquad Z=\sum_i e_iz_i.
\]
From the finite eigen-equations $bz_i=aw_i$ it proves $bZ=aW$.
If $W=1$, then $b\mid a$. A supplied reducedness certificate
$ua+vb=1$ and $b>0$ force $b=1$. The proof is finite integer algebra,
not a call to a spectral theorem.

Equivalently, after primitive scaling the human argument may simply
multiply $z_i=\lambda w_i$ by $e_i$ and add to get
$\lambda=\sum_i e_i z_i\in\Z$. The denominator formulation
matches the explicit integer-coded implementation conveniently.

\code{EigenvalueBound.lean} defines the dot product with one matrix
column and its absolute sum by list recursion. It proves
\[
 \left|\sum_iw_i a_i\right|\le W\sum_i|a_i|
       \quad\hbox{if every }|w_i|\le W.
\]
Given a selected coordinate of absolute value $W>0$ and its eigenvalue
equation, it proves $|\lambda|\le R$ when that column sum is at
most $R$. Constructing the selected coordinate by a maximum scan,
and constructing the primitive vector from rational evaluations, are
explicit remaining host lemmas; neither is hidden in the current
theorem statement.

\subsection{Complete definitions for the rational target}

The target module uses \code{Monomial n} with a rational coefficient
and a length-$n$ natural exponent vector, and
\code{Polynomial n := List (Monomial n)}. Repeated monomials are
allowed. \code{coefficientAt} adds all coefficients with a given
exponent vector. \code{polynomialEq} compares these sums on the
union of the two finite supports; outside that union both sums are zero.
Polynomial addition, subtraction, multiplication and evaluation are
fully defined finite list programs.

A \code{Certificate n s d} contains the $d$ generators, unit
coefficients, unit ideal witnesses, $n$ positive integer denominators,
$n$ integral action matrices and their ideal witnesses. The Lean
representation stores the already-cleared identities
\[
 \Delta_kX_kb_j-\sum_i(A_k)_{ij}b_i
                     =\sum_\ell H_{kj\ell}F_\ell.
\]
The Python representation initially stores rational $M_k$ and then
computes $\Delta_k$; to convert it, multiply its residual witnesses
by $\Delta_k$. Matrix storage order differs explicitly: Python stores
rows; the Lean target stores columns. The mathematical convention is
the same row-vector action $vM$.

The executable target definitions then compute absolute column sums,
the rational coordinate grids, their Cartesian product, the common-zero
test, and its finite filter. The main target is
\begin{verbatim}
def EnumerationGoal : Prop :=
  forall n s d (F : Vector (Polynomial n) s)
      (C : Certificate n s d) (p : Vector Rat n),
    certificateOK F C = true ->
      (p.toList in enumerate F C <->
       zeroTest F p.toList = true)
\end{verbatim}
Again the listing is a legible ASCII rendering; the source is valid
Lean. Arity is enforced by vectors, and the executable zero test
also checks the length of its list input. Invalid certificates produce
an empty result in the target function; completeness is asserted only
when the checker accepts.

\subsection{A noncircular completion plan}

The remaining host proof has the following concrete dependency order.
\begin{enumerate}
\item \textbf{Finite polynomial semantics.} Prove evaluation preserves
addition, scaling and multiplication. Prove that collecting duplicate
monomials preserves evaluation. These finite distributive identities
establish \code{EvaluationTransportGoal}.
\item \textbf{Certificate substitution.} Evaluate each accepted
coefficient identity at the fixed common zero. The ideal-combination
side vanishes by the previous multiplicativity lemma. This proves
\code{CertificateSoundnessGoal} and extracts the unit and action laws.
\item \textbf{Primitive scaling.} Compute a common denominator of the
nonzero evaluation row; compute its gcd and Bezout coefficients. Prove
the resulting integer row is primitive and preserves the eigen-equations.
Instantiate the checked denominator lemma for each coordinate.
\item \textbf{The maximum scan.} Track the first maximizing index of a
nonzero integer row. Prove its entry has positive absolute value and
dominates all entries. Apply the checked column estimate.
\item \textbf{Grid membership.} If $-R_k\le z_k\le R_k$, the index
$z_k+R_k$ belongs to the coordinate range of length $2R_k+1$.
Prove that index returns $z_k/\Delta_k$. Cartesian membership follows
by finite list induction. This establishes \code{GridCompletenessGoal}.
\item \textbf{Filtering.} The checked theorem
\code{enumeration\_of\_gridCompleteness} already derives
\code{EnumerationGoal} from the previous target. Duplicate-freedom
is a separate short consequence of injectivity of each fixed-denominator
coordinate map, followed by product and filter preservation.
\end{enumerate}
The only premise of the final filter theorem is the location theorem,
not the full desired enumeration assertion under a new name. Every
mathematical step needed for that location theorem is proved in the
human argument above.

\subsection{From the host specification to object arithmetic}

The final embedded theorem uses numerical codes, not Lean's types
\code{Rat}, \code{Vector} or arbitrary predicates. Let
$\CertCheck(F,C)$, $Z(F,p)$ and $E(F,C)$ be the specified PR checker,
zero test and list-producing program. With validity and arity tests
included, the target is the first-order sentence
\[
 \forall F,C,p\;\bigl(\CertCheck(F,C)=1\ \land\ \operatorname{Valid}(p)
 \ \longrightarrow\
 [\operatorname{mem}(p,E(F,C))=1\ \leftrightarrow\ Z(F,p)=1]\bigr).
 \tag{E}\label{eq:object-enumeration}
\]
The syntax of the fixed programs is built from zero, successor,
projections, composition and primitive recursion. For each program,
prove its local equations in \code{Prf PRAAx}; then prove its loop
invariant with \code{PRA.qfInd}. The useful invariants are recorded
in Table~\ref{tab:invariants}.

Bounds are computed values of fixed PR functions; they need not be
small. A nested computation whose bound is obtained from an earlier
finite computation remains a PR composition. Proofs keep $F,C,p$
free and quantify only after the finite invariants are established.
This avoids an induction hypothesis containing an unbounded universal
claim about all rational zeros.

\begin{longtable}{P{.22\linewidth}P{.43\linewidth}P{.25\linewidth}}
\caption{Finite invariants needed for the full object derivation.}
\label{tab:invariants}\\
\toprule
Program & Quantifier-free content after bounded Booleanization & Fuel\\
\midrule\endfirsthead
\toprule Program & Invariant content & Fuel\\\midrule\endhead
Extended gcd & Current remainder equations and integer combinations
of the two original inputs. & Sum of input magnitudes plus one.\\[3pt]
Vector Bezout fold & The accumulated gcd is a linear combination of
the processed entries and divides them. & Vector length.\\[3pt]
Coefficient collection & Each processed monomial contributes its
coefficient to the matching exponent slot. & Input list length and
support size.\\[3pt]
Polynomial evaluation & The accumulator is the sum of the values of
the processed monomials. & List length, variable count and exponents.\\[3pt]
Matrix action & Each partial sum agrees with the processed column
terms at the fixed point. & Matrix dimension.\\[3pt]
Maximum scan & The selected index is processed, its value is the
current maximum, and all processed values are at most it. & Vector
length.\\[3pt]
Grid decoding & Each processed digit is in range and the quotient/
remainder reconstruction equation holds. & Number of coordinates.\\[3pt]
Filter & For each processed candidate, output membership is equivalent
to passing the exact equation tests. & Product of coordinate-grid sizes.\\
\bottomrule
\end{longtable}

Once \eqref{eq:object-enumeration} has a \code{Prf PRAAx}
derivation, the existing theorem \code{pra\_iff\_native} transports
it to the native scalar presentation \cite{NativeSource}. That bridge
is a final step. It cannot replace construction of the object derivation.

The printed axiom dependencies also distinguish the layers. The new
object builders and integer eigenvalue lemmas use only the host kernel
principles \code{propext} and \code{Quot.sound}; no
\code{Classical.choice} occurs there. The standard rational-number
implementation used for the host target programs does introduce
\code{Classical.choice} into the host filter theorem's dependencies.
That does not add a rule to \code{PRAAx}. The proposed arithmetic
implementation uses explicit signed-integer and positive-denominator
codes, so its constructive content is not inferred from that host
dependency list. Lean's acceptance of a recursion definition likewise
does not by itself prove primitive recursiveness or reify it into
\code{PR} syntax.

\section{Validation, deliverables, and the next decisive steps}
\label{sec:validation}

\subsection{Exact reference programs}

The Python programs use standard-library arbitrary integers and
\code{fractions.Fraction}; no floating-point decisions enter their
certificate checks or solution lists. The main checker expands the
unit and action residuals as sparse rational polynomials and compares
coefficients. It computes positive common denominators and absolute
column sums, enumerates the mixed-radix coordinate grid, and checks the
original equations at every retained point. Accepted example certificates
are exported as JSON files so the checked data can be inspected or
processed by another implementation.

The featured tests include a four-point elliptic section, rational
coordinates with distinct denominators, a quotient with nilpotents,
an empty rational zero set, a zero-radius grid, a unit expressed as a
linear combination rather than a literal generator, noncommuting
matrices on a redundant spanning family, and a localization that removes
an irrelevant positive-dimensional component. Independent linear-system
and factored-polynomial families compare against solutions known directly
from the input construction. Malformed and altered certificates are
rejected before enumeration. Grid decoding is separately compared with
a Cartesian-product implementation.

The separated-annihilator program also constructs multivariable monomial
rectangles and all action witnesses from the supplied ideal identities.
Its tests cover mixed degrees, degree one, tensor products, repeated roots
and the elliptic section. The norm program independently constructs monic reductions, polynomial
multiplication matrices and their determinants. Its determinant values
are compared with scalar Gaussian elimination at rational sample inputs.
It tests repeated and singular fibers, zero first coordinates, nonzero
constant norms, extraneous norm roots, the final $G=0$ filter, and
nontrivial rational denominators. Its norm-to-module constructor retains
the division quotients and adjugate witnesses. The resulting certificates
are checked by the separate coefficient checker, and its complete lists
are compared with the coordinate-grid enumeration.

These are implementation checks on finite cases. They do not establish
the universal theorems by testing. The universal mathematical arguments
are the proofs in Sections~\ref{sec:finite-observable-modules}
and~\ref{sec:norm-enumeration}; the exact tests detect mistakes in
the concrete programs and examples.

\subsection{Lean and mathematical review}

The standalone Lean build checks the eight modules listed in
Section~\ref{sec:formalization}. The archive records the compiler
version, successful module results, printed axiom dependencies and
the source-provenance check. The example module evaluates the elliptic
certificate independently: the checker returns true, the column bounds
are $[2,1]$, the grid has fifteen candidates, and the result has the
four displayed rational points. An evaluation diagnostic is kept distinct
from a theorem and from an object \ipra\ derivation.

Independent reviews checked the matrix orientation, use of the unit
certificate, denominator divisibility, maximum-coordinate estimate,
constructor residuals, elliptic example, localization identities,
the norm bridge, and the distinction between global and tuple-local
covers. The review records are included. They do not assert that all
earlier repository packets are correct. The PDF is compiled from the
included source and visually inspected after rendering; build and
layout results are recorded separately from mathematical validation.

\subsection{A useful integration order}

The first integration target is the proof-transport core: it is already
compiled and has a small dependency footprint. The next target is the
integer/rational finite arithmetic library described in the formalization
plan. Completing that library supports both enumeration methods and many
existing theorem packets.

For the mathematics, a particularly useful follow-up is a catalogue of
explicit finite sections and their certificates. Each entry should provide
the equations, the norm or action data, an independent identity check,
the returned list, and the precise role of that section in a larger
argument. This is more actionable than adding an abstract field saying
that a desired cover exists. The harder global theorem should then state
exactly which of those certificates it produces and how its coefficients
are bounded.

The current package completes a real interface improvement: finite
algebraic data previously used to count points now supply their
coordinates and complete lists, and a checked logical transformation
preserves local arithmetic proofs through finite computations. The
remaining global arithmetic and analytic obligations are exposed as
specific theorems, with no claim that changing the representation has
already proved them.

\begin{thebibliography}{99}
\small
\bibitem{RepoIndex}
ProveIt contributors. \emph{Primitive recursive arithmetic: directory
overview}. All repository references below are pinned to commit
\path{b0f6f535ffc16324683fef2c263260a88ae03755}.
\href{https://github.com/VladimirReshetnikov/ProveIt/blob/b0f6f535ffc16324683fef2c263260a88ae03755/Logic/PrimitiveRecursiveArithmetic/README.md}{Pinned overview}.

\bibitem{PacketIndex}
ProveIt contributors. \emph{Delivered iPRA proof packets},
\path{Logic/PrimitiveRecursiveArithmetic/docs/proofs/README.md}.
\href{https://github.com/VladimirReshetnikov/ProveIt/blob/b0f6f535ffc16324683fef2c263260a88ae03755/Logic/PrimitiveRecursiveArithmetic/docs/proofs/README.md}{Pinned packet index}.
Subsequent proof-packet paths are relative to this \path{docs/proofs/} directory.

\bibitem{Embedding}
\emph{Mathematics in iPRA through finite observations},
\path{foundations/ipra_embedding_framework/paper/05_diophantine.tex}.
In particular, ``Finite algebra certificates as a general zero-count
interface.''
\href{https://github.com/VladimirReshetnikov/ProveIt/blob/b0f6f535ffc16324683fef2c263260a88ae03755/Logic/PrimitiveRecursiveArithmetic/docs/proofs/foundations/ipra_embedding_framework/paper/05_diophantine.tex}{Pinned section}.

\bibitem{Observation}
\emph{Finite observations and uniform constructions},
\path{foundations/ipra_finite_observation/tex/sections.tex}.
\href{https://github.com/VladimirReshetnikov/ProveIt/blob/b0f6f535ffc16324683fef2c263260a88ae03755/Logic/PrimitiveRecursiveArithmetic/docs/proofs/foundations/ipra_finite_observation/tex/sections.tex}{Pinned finite-sections and cutoff discussion}.

\bibitem{Demand}
\emph{Finite demands, root supports and norm certificates},
\path{foundations/ipra_finite_demand/sections/05_norm_certificates.tex}
and \path{sections/07_audit.tex}.
\href{https://github.com/VladimirReshetnikov/ProveIt/blob/b0f6f535ffc16324683fef2c263260a88ae03755/Logic/PrimitiveRecursiveArithmetic/docs/proofs/foundations/ipra_finite_demand/sections/05_norm_certificates.tex}{Pinned norm section}.

\bibitem{Coherent}
\emph{Coherent finite methods},
\path{foundations/ipra_coherent_methods_2026-11-15/paper/sections/03_refinement.tex},
\path{05_norms.tex}, and \path{07_status.tex}.
\href{https://github.com/VladimirReshetnikov/ProveIt/blob/b0f6f535ffc16324683fef2c263260a88ae03755/Logic/PrimitiveRecursiveArithmetic/docs/proofs/foundations/ipra_coherent_methods_2026-11-15/paper/sections/07_status.tex}{Pinned status and boundary analysis}.

\bibitem{Siegel}
\emph{Siegel's theorem in iPRA: constructive cases and remaining route},
\path{diophantine-number-theory/siegel_ipra/paper/05_general_route.tex}.
\href{https://github.com/VladimirReshetnikov/ProveIt/blob/b0f6f535ffc16324683fef2c263260a88ae03755/Logic/PrimitiveRecursiveArithmetic/docs/proofs/diophantine-number-theory/siegel_ipra/paper/05_general_route.tex}{Pinned finite-cover route}.

\bibitem{Faltings}
\emph{Faltings' theorem in iPRA: scope and constructive special cases},
\path{diophantine-number-theory/faltings_ipra/README.md}.
\href{https://github.com/VladimirReshetnikov/ProveIt/blob/b0f6f535ffc16324683fef2c263260a88ae03755/Logic/PrimitiveRecursiveArithmetic/docs/proofs/diophantine-number-theory/faltings_ipra/README.md}{Pinned scope statement}.

\bibitem{MordellWeil}
\emph{Constructive Mordell--Weil descent},
\path{diophantine-number-theory/mordell_weil_ipra/paper/03_elliptic.tex}.
\href{https://github.com/VladimirReshetnikov/ProveIt/blob/b0f6f535ffc16324683fef2c263260a88ae03755/Logic/PrimitiveRecursiveArithmetic/docs/proofs/diophantine-number-theory/mordell_weil_ipra/paper/03_elliptic.tex}{Pinned elliptic-curve interface}.

\bibitem{PrimeDistribution}
\emph{Prime distribution in iPRA},
\path{analytic-number-theory/prime_distribution_ipra/README.txt}.
\href{https://github.com/VladimirReshetnikov/ProveIt/blob/b0f6f535ffc16324683fef2c263260a88ae03755/Logic/PrimitiveRecursiveArithmetic/docs/proofs/analytic-number-theory/prime_distribution_ipra/README.txt}{Pinned scope statement}.

\bibitem{PRSource}
ProveIt contributors. \emph{Intuitionistic primitive recursive arithmetic},
\path{Logic/PrimitiveRecursiveArithmetic/ConciseIPRA/Lean/ConciseIPRA/PR.lean}.
\href{https://github.com/VladimirReshetnikov/ProveIt/blob/b0f6f535ffc16324683fef2c263260a88ae03755/Logic/PrimitiveRecursiveArithmetic/ConciseIPRA/Lean/ConciseIPRA/PR.lean}{Pinned PR syntax and axioms}.

\bibitem{DeductionSource}
ProveIt contributors. \emph{Scoped syntax and intuitionistic natural deduction},
\path{ConciseIPRA/Lean/ConciseIPRA/Syntax.lean} and
\path{Deduction.lean} within the arithmetic directory.
\href{https://github.com/VladimirReshetnikov/ProveIt/blob/b0f6f535ffc16324683fef2c263260a88ae03755/Logic/PrimitiveRecursiveArithmetic/ConciseIPRA/Lean/ConciseIPRA/Deduction.lean}{Pinned deduction source}.

\bibitem{NativeSource}
ProveIt contributors. \emph{Native bi-interpretation},
\path{ConciseIPRA/Lean/ConciseIPRA/NativeBiInterpretation.lean}.
\href{https://github.com/VladimirReshetnikov/ProveIt/blob/b0f6f535ffc16324683fef2c263260a88ae03755/Logic/PrimitiveRecursiveArithmetic/ConciseIPRA/Lean/ConciseIPRA/NativeBiInterpretation.lean}{Pinned transfer theorem}.

\bibitem{Rouillier}
F. Rouillier. \emph{Solving Zero-Dimensional Systems Through the
Rational Univariate Representation}. Applicable Algebra in Engineering,
Communication and Computing \textbf{9}(5), 433--461 (1999).
\href{https://doi.org/10.1007/s002000050114}{DOI: 10.1007/s002000050114}.
Bibliographic details checked against
\href{https://who.rocq.inria.fr/Fabrice.Rouillier/publications.html}{the author's publication list}.
This is background context; no theorem from it is required by the proofs here.
\end{thebibliography}

\end{document}
